\documentclass[10pt,a4paper]{amsart}
\usepackage[margin=19mm]{geometry}
\usepackage{amsmath,amssymb,mathtools}
\usepackage[scr=rsfs,cal=euler]{mathalfa}
\usepackage{xfrac}
\usepackage[unicode,hidelinks,bookmarks=false]{hyperref}
\newcommand{\ie}{i.e.\ }
\newcommand{\cf}{cf.\ }
\newcommand{\resp}{respectively\ }
\newcommand{\Subsection}{\subsection}
\providecommand{\nobreakdash}{\nobreak\hbox{-}}
\setlength{\emergencystretch}{2em}
\multlinegap0pt
\usepackage{tikz}
\usepackage{mathtools}
\usepackage{amssymb}
\usepackage{dsfont}
\usepackage{bm}
\newtheorem{theorem}{Theorem}
\newtheorem{scope}{Scope}
\def\CC{\mathbb{C}}
\def\NN{\mathbb{N}}
\def\ZZ{\mathbb{Z}}
\newcommand{\dint}{\textup{d}}
\title{\bfseries Random convex chains through the lens\\ of analytic combinatorics}
\author{Florian Besau, Christoph Th\"ale}
\date{Source: arXiv:2510.16793v2 (CC BY 4.0)}
\begin{document}
\maketitle

\noindent\emph{Source and licence.} \bfseries Random convex chains through the lens\\ of analytic combinatorics. Version arXiv:2510.16793v2 (CC BY 4.0), distributed by arXiv under the Creative Commons Attribution 4.0 International licence, http://creativecommons.org/licenses/by/4.0/, as recorded in the arXiv licence metadata for this exact version and shown on the abstract page https://arxiv.org/abs/2510.16793v2. Full licence notice: \texttt{licenses/arxiv-2510.16793-CC-BY-4.0.txt}.\par\smallskip

\noindent\textit{Editorial scope.} Counted unit: the article's theorem thm:saddle (source0.tex lines 1452--1463). The article's complete original proof of it is delivered with it (source0.tex lines 1464--1621, one \texttt{proof} environment of 158 lines). No mathematics is written, repaired or re-derived: the statement and the proof are the article's own lines, in the article's own notation, and no retained line is substituted or re-flowed.  Retained uncounted as the source-local context the proof needs: lines 1189--1199.  One declared adaptation: exactly 3 ancillary strings inside the retained spans are omitted (image-inclusion commands and, where the source repeats a label, the repeated label definitions) and no other line differs from the source; the figure environments, with their placement, captions and labels, are kept, and the package records the omitted strings in fidelity receipts bound to the affected bodies.  No retained line contains a citation, so the wrapper carries no bibliography and no bibliography entry.  The retained lines contain the article's own diagram code, which the wrapper typesets with the standard tikz packages; no image file is delivered and the package declares no asset dependency.  Editorial formatting only, in addition: an AMS \texttt{amsart} preamble that restates the article's own declarations and macros used by the retained lines, namely the environments \texttt{theorem}, \texttt{scope} on separate counters, together with the standard packages those lines need.  The retained environments print in this document as Theorem 1, Scope 1 in order of appearance, whereas the article prints the same items with the numbering of its own full document.  The complete native source of the article as distributed in the arXiv e-print for this exact version is bundled unchanged as \texttt{sources/187363/source0.tex}.
\begin{theorem}[Probability asymptotic for fixed $k$]\label{thm:probab}
    For $k\in\ZZ_{\geq 0}$ we have
    \begin{equation*}
        p_{k+1}^{(n)} = \frac{2^{k+1}}{k!} \frac{(\log n)^{k}}{n} (1+O(1/\log n)),
    \end{equation*}
    and
    \begin{equation}\label{eqn:pnnk}
        p_{n-k}^{(n)} = \frac{n^{3k-1}}{3^k k!} \frac{2^n}{(n!)^2} (1+O(1/n)),
    \end{equation}
    as $n\to\infty$.
\end{theorem}

\begin{theorem}\label{thm:saddle}
    For $c\in(0,\frac{3}{4})$ and a sequence $\{k_n\}_{n\in\NN}$ defined by $k_n:=\lfloor c\log n\rfloor$ we have
    \begin{align*}
        p_{k_n}^{(n)} &= \frac{\sigma(\alpha_*)}{\sqrt{2\pi}} \left(\frac{1}{\alpha_*} + \frac{1}{\alpha_*+1}\right) \frac{K(\alpha_*)}{\Gamma(\alpha_*)} \binom{\alpha_*+1}{2}^{\!\!-k_n} 
            \frac{n^{\alpha_*-1}}{\sqrt{k_n}}\, (1+o(1)),
    \end{align*}
    where 
    \begin{equation*}
        \alpha_* := c-\frac{1}{2} + \frac{1}{2} \sqrt{4c^2+1}>0 \quad \text{and} \quad
        \sigma(\alpha_*) := \left(\frac{1}{\alpha_*^2} + \frac{1}{(\alpha_*+1)^2}\right)^{-1/2}>0.
    \end{equation*}
\end{theorem}

\begin{proof}
    As in the beginning of the proof of Theorem \ref{thm:probab} and using the Cauchy integral formula, we derive
    \begin{align*}
        p_{k_n}^{(n)} &= [z^{k_n}] G_n(z) \\
        &= [\alpha^{k_n}] \frac{1+2\alpha}{1+\alpha}\frac{K(\alpha)}{\Gamma(\alpha)} \left(\frac{2}{1+\alpha}\right)^{k_n} n^{\alpha-1} (1+O(1/n))\\
        &= \frac{1}{2\pi i} \oint \left(\frac{1}{\alpha}+\frac{1}{1+\alpha}\right)\frac{K(\alpha)}{\Gamma(\alpha)} \left(\frac{2}{\alpha(1+\alpha)}\right)^{k_n} 
             n^{\alpha-1} \,\dint \alpha\, (1+O(1/n))\\
        &= \frac{1}{2\pi i} \oint H(\alpha) \binom{\alpha+1}{2}^{-k_n}
            n^{\alpha-1} \, \dint\alpha\, (1+O(1/n)),
    \end{align*}
    where we set
    \begin{equation*}
        H(\alpha) := \left(\frac{1}{\alpha}+\frac{1}{\alpha+1}\right) \frac{K(\alpha)}{\Gamma(\alpha)} = \frac{1}{(1+\alpha)^2} \frac{\Gamma(2+2\alpha)}{\Gamma(1+\alpha)^3}.
    \end{equation*}
    The complex contour integral is over any simple loop $\alpha(t)$ in $\CC$ around $0$ with $\mathrm{Im}\, \alpha>-1$, since the integrand has a pole at $\alpha=-1$ and is analytic for $\mathrm{Im}\, \alpha>-1$ apart from the pole at $\alpha=0$ of order $k_n$.

By Cauchy's Theorem we may choose to consider the contour integral over a rectangle with vertices $-\sigma\pm i T$ and $b\pm T$ for some $\sigma\in (0,1)$ and large $T>0$.
We will show that the main contribution of the contour is given by the vertical contour $b+iT$, see Figure \ref{fig:graph_In}.

\begin{figure}[ht]
    \centering
    \begin{tikzpicture}
        \begin{scope}
            \clip (-2.5,-2.5) rectangle (2.5,2.5);
            \node at (0,0) {};
            %\draw[color=black!80, ->] (-0.65,-1.33) to (-2,0) node[left] {$it$};
            %\draw[color=black!80, ->] (-0.65,-1.33) to (2,-0.63) node[right] {$s$};
            %\draw (-2.5,-2.5) rectangle (2.5,2.5);
        \end{scope}
        
        \begin{scope}[xshift=6cm]
            \clip (-2.5,-2.5) rectangle (2.5,2.5);
            \node at (0,0) {};
            %\draw (-2.5,-2.5) rectangle (2.5,2.5);
        \end{scope}
        
        \begin{scope}[xshift=12cm]
            \clip (-2.5,-2.5) rectangle (2.5,2.5);
            \node at (0.2,0) {};
            %\draw (-2.5,-2.5) rectangle (2.5,2.5);
        \end{scope}
    \end{tikzpicture}
    \caption{Graph of $|I_n(s+it)|$ for $n=10^6$ and $c=2/3$ (left), $c=3/4$ (middle), and $c=0.8$ (right). The green curve is the section $s\equiv \tilde{\alpha}_*$ and the red curve is the section $s\equiv \alpha_*$.}
    \label{fig:graph_In}
    % Link to script:
    %https://sagecell.sagemath.org/?z=eJx9Uc1uozAQviPxDkg9MDaGYNNDVck9RqqqPkGUVA5xCA2xqXFbHr9jlmS726iWrPHw_WmYD-UgVWybkjgykpfTiaNalsVdHCmeyKTO-UJkeOnw5jzc0nqDLfKVOKP5FfRo5L6z1kFNO9uAIfitd63xcDT4fEKt2g7AF8CzkWwEbdTppEDQMRNk8aeZkCqhidnAmHNCAUY0p-NCkA3kx8k0jnZ6nywBpyD3cZTgcdqHcOXhqRjeMWWUKnuk20APeLsPlIdqps-Sd2eSqrg4hBbLJYDDd_-AYibH0AtDXGOImXGefhKVf__GRCnJzLnpO-urHSzZSjHF87LgWDIsa7baspwzfNS2s06m1inT6LC5vvPPWv4jxZ0E4SwTTKxZwF96i6GDXOGqGd71WZ3JXjl10t619cvstArjseV3jzna6V3K_KGtj0YPg6x-dxHBRfx0aZzW5n-f2akYDvYTemdfde1ba8K4_mAbp3qkp-QLdw6_Zg==&lang=sage&interacts=eJyLjgUAARUAuQ==
\end{figure}

\medskip 
\emph{Estimating the horizontal contours:} 
By Stirling's formula we may estimate for all $s>0$,
\begin{align*}
    |\Gamma(s+iT)| &= e^{-\frac{\pi}{2} |T| (1+o(1))} \quad \text{as $|T| \to \infty$}
\end{align*}
which yields for all $s\in [-\sigma,b]\subset (-1,+\infty)$
\begin{equation*}
    |H(s+iT)| \leq e^{C_1|T|},
\end{equation*}
for some $C_1=C_1(\sigma,b)>0$ and $|T|$ large enough.
Furthermore,
\begin{equation*}
    \left|\binom{1+s+iT}{2}^{-k_n}\right| 
    \leq n^{-C_2|T|},
\end{equation*}
for some $C_2=C_2(\sigma,b,c)>0$, and
\begin{equation*}
    |n^{s+iT-1}| = n^{s-1}\leq n^{b-1}.
\end{equation*}
Writing
\begin{equation*}
    I_n(\alpha) := H(\alpha)\, \binom{\alpha+1}{2}^{-k_n}\, n^{\alpha-1}, 
    %= H(\alpha) n^{\alpha - 1 - c\log\binom{\alpha+1}{2} + O(1/\log n)},
\end{equation*}
we conclude
\begin{equation*}
    \left|\int_{-\sigma}^b I_n(s+iT)\, \dint s \right| \leq (b+\sigma) n^{C_1\frac{|T|}{\log n}-C_2|T|+b-1}.
\end{equation*}
So for $T=T_n:=n^2$ we find
\begin{equation*}
    \left|\frac{1}{2\pi i} \int_{-\sigma}^b I_n(s+iT_n)\, \dint s\, (1+O(1/n)) \right| \leq n^{-C_3 n^2},
\end{equation*}
for some $C_3=C_3(c,b,\sigma)>0$.
Thus, the horizontal parts of our rectangle will be asymptotically small.
\medskip

\emph{Estimating the vertical contours:}
Next, we want to estimate the vertical parts
\begin{equation*}
    \frac{1}{2\pi i}\int_{\mathrm{Re}\, \alpha=s_0, |\mathrm{Im}\,\alpha|\leq T_n} I_n(\alpha)\, \dint \alpha\, (1+O(1/n)),
\end{equation*}
for suitable fixed $s_0\in (-1,+\infty)$.
For this, we write
    \begin{align*}
        I_n(\alpha)&= n^{S_n(\alpha)}
    \end{align*}
where we set
\begin{equation*}
    S_n(\alpha) := \alpha - 1 + \frac{\log H(\alpha)}{\log n}-(\log \alpha) \frac{k_n}{\log n} - \log\left(\frac{1+\alpha}{2}\right)\frac{k_n}{\log n}.
\end{equation*}
Then
\begin{equation*}
    S_n'(\alpha) = 1 + \frac{H'(\alpha)}{H(\alpha)\log n} - \frac{1}{\alpha} \frac{k_n}{\log n} - \frac{1}{1+\alpha} \frac{k_n}{\log n} + O\left(\frac{1}{\log n}\right).
\end{equation*}
We set
\begin{equation*}
    b=\alpha_* := c-\frac{1}{2}+\frac{1}{2}\sqrt{4c^2+1} >0 \quad \text{and}\quad
    -\sigma=\tilde{\alpha}_* := c-\frac{1}{2}-\frac{1}{2}\sqrt{4c^2+1} \in (-1,-1/2),
\end{equation*}
and conclude that for $\alpha^\dag\in\{\alpha_*,\tilde{\alpha}_*\}$
\begin{equation*}
    \lim_{n\to\infty} S_n'(\alpha^\dag) = 1-c\left(\frac{1}{\alpha^\dag} +\frac{1}{1+\alpha^\dag}\right)=0,
\end{equation*}
which in turn yields
    \begin{equation*}
        S_n'(\alpha^\dag) = \frac{H'(\alpha^\dag)}{H(\alpha^\dag)\log n} + \left(\frac{1}{\alpha^\dag} +\frac{1}{\alpha^\dag+1}\right) \frac{c\log n-\lfloor c\log n\rfloor}{\log n} 
             + O\left(\frac{1}{\log n}\right) 
            = O\left(\frac{1}{\log n}\right).
    \end{equation*}
    Furthermore,
    \begin{equation*}
        S''_n(\alpha^\dag) = \left(\frac{1}{(\alpha^\dag)^2} + \frac{1}{(1+\alpha^\dag)^2}\right) \frac{k_n}{\log n} + O\left(\frac{1}{\log n}\right) 
        = \frac{c}{\sigma(\alpha^\dag)^2} \left(1+ O\left(\frac{1}{\log n}\right)\right).
    \end{equation*}
    It follows that
    \begin{equation*}
        S_n(\alpha) = S_n(\alpha^\dag) + \frac{1}{2} S''_n(\alpha^\dag) (\alpha-\alpha^\dag)^2 + O\left(\frac{\alpha-\alpha^\dag}{\log n}\right) + O((\alpha-\alpha^\dag)^3)
    \end{equation*}
    So on the vertical contours $\mathrm{Re}\,\alpha=\alpha^\dag$ we have
    \begin{equation*}
        S_n(\alpha^\dag+it) = S_n(\alpha^\dag) - \frac{t^2}{2} S''_n(\alpha^\dag) + iO\left(\frac{t}{\log n}+t^3\right),
    \end{equation*}
    and therefore
    \begin{align*}
            &\frac{1}{2\pi} \int_{-T_n}^{T_n} I_n(\alpha^\dag+it) \, \dint t\, (1+O(1/n))\\
            &\qquad=\frac{1}{2\pi} \int_{-T_n}^{T_n} n^{S_n(\alpha^\dag)-\frac{t^2}{2} S''_n(\alpha^\dag)} e^{i O(t+(\log n)t^3)} \, \dint t\, (1+O(1/n))\\
            &\qquad= \frac{n^{S(\alpha^\dag)}}{2\pi \sqrt{(\log n)S''_n(\alpha^\dag)}} \int_{-T_n\sqrt{(\log n)S''_n(\alpha^\dag)}}^{T_n\sqrt{(\log n)S''_n(\alpha^\dag)}} e^{-s^2/2} e^{i O\left(\frac{s}{\sqrt{\log n}}\right)} \,\dint s\, (1+O(1/\sqrt{\log n}))\\
            &\qquad= \frac{\sigma(\alpha^\dag)}{2\pi} \frac{n^{S(\alpha^\dag)}}{\sqrt{k_n}} \int_{-\infty}^{+\infty} e^{-s^2/2} \left(1+O\left(\frac{s}{\sqrt{\log n}}\right)\right)\,\dint s\, (1+O(1/\sqrt{\log n}))\\
            &\qquad= \frac{\sigma(\alpha^\dag)}{\sqrt{2\pi}} \frac{n^{ S_n(\alpha^\dag)}}{\sqrt{k_n}} (1+O(1/\sqrt{\log n})).
    \end{align*}
    For $\alpha^\dag=\tilde{\alpha}_*$ we further estimate
    \begin{equation*}
        \left|\frac{1}{2\pi} \int_{-T_n}^{T_n} I_n(\tilde{\alpha}_*+it) \, \dint t\, (1+O(1/n)) \right|
        \leq \frac{C_4}{\sqrt{k_n}} n^{\tilde{\alpha}_*-1-c \log|\binom{\tilde{\alpha}_*+1}{2}|},
    \end{equation*}
    and conclude
    \begin{align*}
        p_{k_n}^{(n)} 
        &= \frac{1}{2\pi} \int_{-T_n}^{T_n} I_n(\alpha_*+it) \, \dint t (1+O(1/n)) 
            +  \frac{C_4}{\sqrt{k_n}} n^{\tilde{\alpha}_*-1-c \log|\binom{\tilde{\alpha}_*+1}{2}|} + n^{-C_3n^2}\\
        &=\frac{\sigma(\alpha_*)}{\sqrt{2\pi}} \frac{H(\alpha_*)}{\sqrt{k_n}}n^{\alpha_*-1-c\log\binom{\alpha_*+1}{2}} (1+O(1/\sqrt{\log n}))
            +  \frac{C_4}{\sqrt{k_n}} n^{\tilde{\alpha}_*-1-c \log|\binom{\tilde{\alpha}_*+1}{2}|} + n^{-C_3n^2}.
    \end{align*}
    This completes the proof, since we can check that 
    \begin{equation*}
        -I(c) = \alpha_*-1-c\log\binom{\alpha_*+1}{2} > \tilde{\alpha}_*-1-c \log\left|\binom{\tilde{\alpha}_*+1}{2}\right|,
    \end{equation*}
    for $c\in (0,c_0)$, where we choose $c_0:=\frac{3}{4}$ somewhat arbitrarily.
\end{proof}

\end{document}
